% This document was generated by an LLM.
% Content fragment; compile intro-analysis-self-check.tex.
% Coverage: Hunter (2012), Chapter 5, printed pp. 57--72.
% Source: https://www.math.ucdavis.edu/~hunter/m125a/intro_analysis.pdf
% Corrections and clarifications:
% - Bounds for a sequence of functions distinguish each fixed input from
%   a common bound over both the domain and all indices (Example 5.4).
% - In Example 5.9, choose an integer N and treat b=0 separately.
% - In Example 5.23, sum rho^n=1/(1-rho), not a number less than 1.
% - Example 5.24 is indexed from n=1, unlike its figure caption.
% - In Theorem 5.33, homogeneity ends with |k| times the norm, not k.
% - For the sup-norm questions, assume nonempty domains so that no
%   convention for the supremum over an empty domain is needed.
% - Completeness of C(K) uses continuity of uniform limits as well as
%   the uniform Cauchy criterion (Theorem 5.36).

\section{Sequences and Series of Functions}

\subsection{Pointwise convergence}
\textit{Source: \S5.1, pp.~57--59.}

\begin{enumerate}
  % Definition 5.1.
  \item Define pointwise convergence of \(f_n:A\to\R\) to
    \(f:A\to\R\). Write the definition with all quantifiers.
    On which quantities may the index threshold depend?

  % Example 5.2.
  \item Find the pointwise limit of \(f_n(x)=n/(nx+1)\) on \((0,1)\).
    Show that each \(f_n\) is bounded on this domain and determine
    whether its limit is bounded. What preservation claim does this disprove?

  % Example 5.3.
  \item Find the pointwise limit of \(f_n(x)=x^n\) on \([0,1]\).
    Is the limit continuous? What does this show about pointwise
    limits of continuous functions?

  % Example 5.4; explicit formula makes the question independent of the PDF.
  \item On \([0,1]\), define
    \[
      f_n(x)=
      \begin{cases}
        2n^2x,&0\leq x\leq 1/(2n),\\
        2n^2(1/n-x),&1/(2n)<x<1/n,\\
        0,&1/n\leq x\leq1.
      \end{cases}
    \]
    Find its pointwise limit, its maximum, and the location of that
    maximum. Why is \((f_n(x))\) bounded for each fixed \(x\),
    even though there is no common bound over all \(x\) and \(n\)?

  % Example 5.5.
  \item Find the pointwise limit of \(f_n(x)=\sin(nx)/n\) on \(\R\).
    Does the sequence of derivatives converge pointwise on \(\R\)?
    Exhibit a point where it fails and explain what this says about
    interchanging differentiation and a pointwise limit.

  % Example 5.6.
  \item Find the pointwise limit on \(\R\) of
    \[
      f_n(x)=\frac{x^2}{\sqrt{x^2+1/n}}.
    \]
    Explain why every approximating function is differentiable but
    the limit need not be differentiable everywhere.

  % Example 5.7; the underlying exponential limit is quoted, not proved.
  \item State the pointwise limit of \((1+x/n)^n\) on \(\R\),
    using the exponential limit formula quoted in Example 5.7.
    Does identifying this pointwise limit by itself establish uniform
    convergence on \(\R\)?
\end{enumerate}

\subsection{Uniform convergence}
\textit{Source: \S5.2, p.~59.}

\begin{enumerate}[resume]
  % Definition 5.8.
  \item Define uniform convergence on \(A\), and compare its
    quantifier order with pointwise convergence. Which dependence
    of the threshold \(N\) is forbidden in the uniform definition?

  % Discussion following Definition 5.8.
  \item Prove that uniform convergence implies pointwise convergence
    to the same limit. Write the negation of uniform convergence
    to a specified function, including all quantifiers.

  % Example 5.9.
  \item Prove that \(x^n\) does not converge uniformly on \([0,1]\).
    Does it converge uniformly on \([0,b]\) for \(0\leq b<1\)?
    Give a suitable integer threshold for a specified error tolerance,
    treating \(b=0\) separately.

  % Example 5.10; refer to the fully specified sequence above.
  \item Show that the triangular functions defined in \S5.1 do not
    converge uniformly to their pointwise limit. Which moving input
    exposes the failure, and why is checking fixed inputs insufficient?

  % Example 5.11.
  \item Prove that \(\sin(nx)/n\) converges uniformly on \(\R\)
    and give an explicit uniform error bound.
    Does this stronger convergence fix the failure of derivative
    convergence found in \S5.1?
\end{enumerate}

\subsection{Cauchy condition for uniform convergence}
\textit{Source: \S5.3, pp.~60--61.}

\begin{enumerate}[resume]
  % Definition 5.12.
  \item Define a uniformly Cauchy sequence of functions on \(A\).
    How does this differ from requiring a Cauchy sequence of real
    values separately at every fixed input?

  % Theorem 5.13, forward direction.
  \item State the uniform Cauchy criterion (Theorem 5.13) and prove
    that a uniformly convergent sequence is uniformly Cauchy.

  % Theorem 5.13, converse direction.
  \item Prove the converse of Theorem 5.13. First construct the
    pointwise limit, then show uniform convergence to it.
    Where is completeness of \(\R\) used? If you pass to the
    limit in a strict Cauchy bound, why might the resulting inequality
    be non-strict, and how do you still obtain the desired error bound?
\end{enumerate}

\subsection{Properties of uniform convergence}
\textit{Source: \S5.4, pp.~61--65.}

\subsubsection*{Boundedness and continuity}
\begin{enumerate}[resume]
  % Theorem 5.14.
  \item State and prove the boundedness theorem for uniform limits
    (Theorem 5.14). Is a common bound for all approximating functions
    assumed, or only boundedness of each individual function?
    Why does choosing one sufficiently late function suffice in the proof?

  % Discussion following Theorem 5.14.
  \item If \(|f_n(x)|\leq M\) for every \(n\) and every \(x\in A\),
    must a pointwise limit be bounded by \(M\)?
    Distinguish this hypothesis from individual boundedness of each \(f_n\).

  % Example 5.15.
  \item For \(f_n(x)=n/(nx+1)\), explain why convergence on \((0,1)\)
    is not uniform. Prove uniform convergence on \([a,1)\) for
    \(0<a<1\), giving an error bound in terms of \(a\) and \(n\).

  % Theorem 5.16.
  \item State the uniform-limit theorem for continuous functions
    (Theorem 5.16). Is compactness of the domain a hypothesis?

  % Proof of Theorem 5.16.
  \item Prove Theorem 5.16 using an \(\epsilon/3\) argument.
    Write the three-term estimate for \(|f(x)-f(c)|\).
    Why do you fix an approximating index before choosing the
    continuity radius?

  % Discussion following Theorem 5.16.
  \item Express the continuity theorem as an exchange of the limits
    \(n\to\infty\) and \(x\to c\), at an accumulation point of
    the domain. Use \(x^n\) at \(c=1\) to show why pointwise
    convergence alone does not justify this exchange.

  % Final paragraph of Section 5.4.2.
  \item If continuous functions converge pointwise to a discontinuous
    function, what follows about uniform convergence?
    Is continuity of the pointwise limit sufficient for uniform
    convergence? Give a counterexample from this chapter.
\end{enumerate}

\subsubsection*{Differentiability}
\begin{enumerate}[resume]
  % Example 5.17.
  \item For \(f_n(x)=x/(1+nx^2)\) on \(\R\), find its uniform
    limit and an explicit uniform error bound.
    Compute \(f_n'\) and its pointwise limit. Does that limit equal
    the derivative of the limiting function? Is convergence of
    the derivatives uniform?

  % Theorem 5.18.
  \item State the theorem justifying differentiation of a sequence
    of functions (Theorem 5.18). Which convergence is pointwise
    and which must be uniform? Are continuity of the derivatives
    or uniform convergence of the original functions assumed?

  % Key mean-value estimate in the proof of Theorem 5.18.
  \item In the proof of Theorem 5.18, apply the mean value theorem
    to \(f_m-f_n\) between distinct \(c\) and \(x\).
    Use uniform Cauchy bounds for the derivatives and then let
    \(m\to\infty\) to control
    \[
      \left|\frac{f(x)-f(c)}{x-c}
        -\frac{f_n(x)-f_n(c)}{x-c}\right|.
    \]
    Why is the resulting bound uniform in the choice of \(x\ne c\)?

  % Remaining epsilon/3 argument in Theorem 5.18.
  \item Complete the proof of Theorem 5.18 by estimating the difference
    between the difference quotient of \(f\) and \(g(c)\) with
    three error terms. Explain the order of the choices of \(n\)
    and the input radius, and why no continuity assumption on \(g\)
    is needed in this proof.

  % FTC discussion on p. 65.
  \item If the derivatives in Theorem 5.18 are also continuous,
    outline the alternative proof using the fundamental theorem of
    calculus on a compact subinterval. Which passage of a limit
    through an integral does uniform convergence justify?
\end{enumerate}

\subsection{Series}
\textit{Source: \S5.5, pp.~65--66.}

\begin{enumerate}[resume]
  % Definition 5.19.
  \item Define the partial sums of a series of functions, and define
    pointwise and uniform convergence of the series in terms of them.
    Is uniform convergence of the individual terms to zero the same
    as uniform convergence of the series?

  % Example 5.20.
  \item Compute the finite partial sums of the geometric series
    \(\sum_{k=0}^{\infty}x^k\). Determine all real inputs where
    it converges, its sum there, and its behavior at the two endpoints.

  % Uniformity discussion in Example 5.20.
  \item Prove that the geometric series is not uniformly convergent
    on \((-1,1)\), but is uniformly convergent on \([-\rho,\rho]\)
    for every \(0\leq\rho<1\). Give a uniform bound for the
    remainder after the \(n\)th partial sum.

  % Theorem 5.21 and its proof.
  \item State and prove the uniform Cauchy criterion for series
    (Theorem 5.21). Which finite tails must be small simultaneously,
    and how does this follow from the sequence version?
\end{enumerate}

\subsection{The Weierstrass \texorpdfstring{\(M\)}{M}-test}
\textit{Source: \S5.6, pp.~67--69.}

\begin{enumerate}[resume]
  % Theorem 5.22.
  \item State the Weierstrass \(M\)-test (Theorem 5.22), specifying
    the domain-independent bounds and the numerical convergence
    hypothesis. Must the sum of the function series be known in advance?

  % Proof of Theorem 5.22.
  \item Prove the \(M\)-test using the uniform Cauchy criterion
    and a triangle-inequality estimate on finite tails.

  % Example 5.23; correct the printed bound on the majorant sum.
  \item Apply the \(M\)-test to \(\sum_{n=0}^{\infty}x^n\) on
    \([-\rho,\rho]\), where \(0\leq\rho<1\).
    Identify the majorants and compute their sum.

  % Example 5.24; nowhere differentiability is stated without proof in source.
  \Needspace{10\baselineskip}
  \item Consider
    \[
      f(x)=\sum_{n=1}^{\infty}2^{-n}\cos(3^n x).
    \]
    Prove uniform convergence on \(\R\) and continuity of \(f\).
    Write the formal term-by-term derivative series.
    What striking differentiability property of \(f\) is reported
    in Example 5.24, without proof?

  % Absolute convergence discussion following Example 5.24.
  \item Define absolute convergence of a series of functions at each
    point of a domain. Why does the \(M\)-test imply absolute
    convergence as well as uniform convergence?
    Explain why the two notions of convergence are different.

  % Example 5.25.
  \item Give a series of constant functions that converges uniformly
    on \(\R\) but is not absolutely convergent at any point.
    Verify both claims and explain why the \(M\)-test cannot apply
    to this series.

  % Contrast with Example 5.20.
  \item Give a series that converges absolutely at every point of
    its domain but does not converge uniformly there.
    Why is pointwise absolute convergence not a sufficient replacement
    for the summable-majorant hypothesis in the \(M\)-test?
\end{enumerate}

\subsection{The sup-norm}
\textit{Source: \S5.7, pp.~69--70. Assume \(A\ne\varnothing\).}

\begin{enumerate}[resume]
  % Definition 5.26.
  \item Define the sup-norm of a function on \(A\).
    When is it finite? Must the supremum be attained, and is the
    expression a finite-valued norm on all real-valued functions?

  % Example 5.27.
  \item Compute the sup-norms on \((0,1)\) of \(x^2\), \(x^2-x\),
    and \(1/x\). For each finite supremum, decide whether it is attained.

  % Definition 5.28.
  \item Characterize uniform convergence using \(\|f_n-f\|\).
    Explain why this is equivalent to the quantified definition,
    even though a pointwise strict bound may give only a non-strict
    bound after taking a supremum.

  % Definition 5.29.
  \item Formulate the uniform Cauchy condition using the sup-norm.
    How does the notation parallel the Cauchy condition for real numbers?
\end{enumerate}

\subsection{Spaces of continuous functions}
\textit{Source: \S5.8, pp.~70--72. Assume \(K\) is nonempty and compact.}

\begin{enumerate}[resume]
  % Definition 5.30 and Theorem 5.31.
  \item Define \(C(K)\). State its vector-space operations and
    explain why they preserve membership in \(C(K)\).
    What are its zero vector and additive inverses?

  % Definition 5.32.
  \item Define a norm on a real vector space and a normed vector
    space. State all norm axioms, including finiteness, and explain
    how a norm determines a distance between vectors.

  % Theorem 5.33 and its proof.
  \item Prove that the sup-norm makes \(C(K)\) a normed vector
    space (Theorem 5.33). Where is compactness needed?
    In the homogeneity identity, what happens when the scalar is negative?

  % Definition 5.34.
  \item Define convergence and the Cauchy property in a normed
    vector space. What do they mean specifically for \(C(K)\)
    equipped with the sup-norm?

  % Definition 5.35.
  \item Define completeness of a normed vector space and a Banach
    space. Must the limit of a Cauchy sequence belong to the same space?

  % Theorem 5.36; make the continuity step explicit.
  \item State and prove completeness of \(C(K)\) in the sup-norm
    (Theorem 5.36). Use both the uniform Cauchy criterion and the
    continuity theorem for uniform limits. Why is existence of a
    real-valued pointwise limit alone insufficient for this conclusion?
\end{enumerate}
